% =====================================================================
% Section: the influence function of the MCD classifier nuisance IF(q_r)
% for logistic regression (classical and ridge), the MLP, and CatBoost.
%
% Organizing principle: a single closed form covers all three, because
% the MLP and CatBoost are handled through a frozen representation and a
% penalized logistic head. The generalized propensity score never
% appears. The Riesz route is kept only as a fallback remark.
%
% Self-contained preamble so it compiles alone; drop-in compatible with
% the companion paper (same notation and equation names).
% =====================================================================

\documentclass[11pt]{article}
\usepackage[T1]{fontenc}
\usepackage[margin=1in]{geometry}
\usepackage{amsmath,amssymb,amsthm}
\usepackage{bm}
\usepackage{mathtools}
\usepackage{hyperref}
\usepackage[round]{natbib}
\usepackage{parskip}
\usepackage{xcolor}
\usepackage[most]{tcolorbox}
\usepackage{booktabs}

\newtheorem{proposition}{Proposition}
\newtheorem{lemma}{Lemma}
\newtheorem{corollary}{Corollary}
\newtheorem*{remark}{Remark}
\newtheorem*{assumption}{Assumption}

\newcommand{\Eb}{\mathbb{E}}
\newcommand{\Tset}{\mathcal{T}}
\newcommand{\Xset}{\mathcal{X}}
\newcommand{\IF}{\mathrm{IF}}
\newcommand{\Hb}{\mathcal{H}}

\newtcolorbox{keyresult}[1][]{
  colback=blue!4, colframe=blue!45!black, boxrule=0.4pt,
  arc=1mm, left=6pt, right=6pt, top=4pt, bottom=4pt, #1
}

\begin{document}

\section{The influence function of the classifier nuisance}
\label{sec:if-qr}

The MCD-native correction of the interventional density is
\begin{equation}
\begin{aligned}
\IF\big(o;p_t(y),F\big)
&= f_Y(y)\,c\left\{ g(t,X_o,y) - \Eb_X[g(t,X,y)]
+ \Eb_X\!\left[\frac{\IF\big(o;q_r(t,X,y),F\big)}{\big(1-q_r(t,X,y)\big)^2}\right] \right\} \\
&\quad + m(t,y)\big(K_{h_Y}(Y_o-y)-f_Y(y)\big),
\end{aligned}
\label{eq:if-mcd-native}
\end{equation}
with $g=q_r/(1-q_r)$, $c=(1-r)/r$ and $m(t,y)=c\,\Eb_X[g(t,X,y)]$. Every term is
elementary except one: the influence function of the classifier itself,
$\IF(o;q_r(t,x,y),F)$. This section computes it for the three binary classifiers we
use, and the outcome is a single closed form covering all three.

The point of the whole construction is to avoid the generalized propensity score
$\pi(t\mid x)$, a conditional density in the multidimensional treatment coordinate,
which the classical correction requires and which is precisely what MCD was designed
to eliminate. Nothing below reintroduces it: the correction is expressed through the
classifier's own parameters, and the only expectations to be estimated are ordinary
sample averages over the augmented data.

We work with the augmented sample $W=(T,X,Y,Z)$ of the MCD construction: with
probability $r$ the triple $(T,X,Y)$ is real and labeled $Z=1$; with probability
$1-r$ it is contrastive, $(T,X,Y')$ with $Y'\sim f_Y$ drawn independently of
$(T,X)$, labeled $Z=0$. Write $P_1(F)=F_{TXY}$ for the real law, $P_0(F)=F_{TX}\otimes
F_Y$ for the contrastive law, and $\Eb_{P_r}=r\,\Eb_{P_1}+(1-r)\,\Eb_{P_0}$ for
expectation under the augmented mixture. The single structural fact driving
everything below is that $P_0(F)$ is itself a functional of $F$, being built from the
marginals of $F$: perturbing $F$ moves the contrastive class as well as the real one.

\subsection{What the classifier must satisfy}
\label{sec:if-qr-regularity}

The object we need, $\IF(o;q_r,F)$, is a derivative: it measures how far the fitted
function $q_r$ moves when the data-generating law $F$ is perturbed infinitesimally
toward a point mass at $o$, along the path $F_\varepsilon=(1-\varepsilon)F+\varepsilon\delta_o$.
For this derivative to exist, the training map $F\mapsto q_r(\cdot;F)$ must be smooth:
a small perturbation of $F$ must displace $q_r$ by a small amount, not make it jump.

A classifier defined as the minimizer of a smooth penalized criterion has this
property, and for a concrete reason. Its parameter $\beta$ is characterized by an
equation, the first-order condition of the criterion. Perturbing $F$ displaces that
equation slightly, hence displaces its solution slightly, and the implicit function
theorem converts the displacement of the equation into the displacement of the
solution. This is the entire mechanism behind every sandwich formula, and it is the
mechanism we use.

A classifier that is not defined this way does not qualify, and two of our three
families do not. A tree selects split thresholds; under a perturbation of $F$ a
threshold eventually flips from one value to another and the fitted function jumps,
so no derivative exists. Boosting compounds the problem: the ensemble is built
greedily, one tree at a time, each correcting the residual of the previous ones, so
there is no single criterion to differentiate at all, only a sequential recipe. A
multilayer perceptron fails for different reasons: its training criterion is
nonconvex, so the minimizer is not a single-valued function of $F$; its Hessian at a
trained network is typically singular; and training is stopped early rather than run
to a stationary point, so the first-order condition does not hold exactly.
\citet{basu2021} document that influence estimates computed as if these conditions
held degrade sharply with depth, width and scale, and \citet{bae2022} show that what
the damped implementation actually approximates is a proximal response around the
trained weights, not the Gateaux derivative of a statistical functional. Since the
bias cancellation in~\eqref{eq:if-mcd-native} is an identity about the true
derivative, an object that only loosely tracks it delivers none of the cancellation.

Section~\ref{sec:if-qr-frozen} resolves this by construction rather than by
approximation: the MLP and CatBoost are split into a representation, which is frozen,
and a logistic head, which is refitted and is a smooth penalized M-estimator by
design. The master formula of Section~\ref{sec:if-qr-master} then applies to all three
classifiers without modification.

\subsection{The master formula}
\label{sec:if-qr-master}

Let the classifier be a penalized empirical risk minimizer over a model
$\{q_{r,\beta}:\beta\in\mathbb{R}^k\}$, with loss $\ell(z,q)$ smooth in $q$ and
penalty $\lambda\mathcal{J}(\beta)$. The population criterion and its minimizer are
\begin{equation}
M_\lambda(\beta;F)
= r\,\Eb_{P_1}\big[\ell(1,q_{r,\beta})\big]
+ (1-r)\,\Eb_{P_0}\big[\ell(0,q_{r,\beta})\big]
+ \lambda\,\mathcal{J}(\beta),
\qquad
\beta_\lambda(F)=\arg\min_\beta M_\lambda(\beta;F).
\label{eq:pop-criterion}
\end{equation}
Write the per-label gradients and the penalized population Hessian
\[
u_\beta=\nabla_\beta\ell(1,q_{r,\beta}),\qquad
v_\beta=\nabla_\beta\ell(0,q_{r,\beta}),\qquad
\Hb_\lambda= r\,\Eb_{P_1}[\nabla_\beta u_\beta]+(1-r)\,\Eb_{P_0}[\nabla_\beta v_\beta]+\lambda\nabla^2\mathcal{J}(\beta).
\]

\begin{assumption}[Smooth penalized classifier]
\label{as:smooth}
$\beta_\lambda(F)$ exists, is unique and lies in the interior of the parameter set;
$\ell(z,\cdot)$ and $\mathcal{J}$ are twice differentiable; $\Hb_\lambda$ is
nonsingular at $\beta_\lambda(F)$; and $q_{r,\beta}$ is bounded away from $0$ and $1$
on the support of interest, so that $g$ and $1/(1-q_r)$ are finite.
\end{assumption}

\begin{proposition}[Three-channel influence function]
\label{prop:master}
Under Assumption~\ref{as:smooth}, with all functions evaluated at
$\beta=\beta_\lambda(F)$ and $o=(y_o,t_o,x_o)$,
\begin{equation}
\IF\big(o;\beta_\lambda,F\big) = -\,\Hb_\lambda^{-1}\big(D_1+D_0^{TX}+D_0^{Y}\big),
\label{eq:master}
\end{equation}
\begin{equation}
\begin{aligned}
D_1 &= r\big( u_\beta(t_o,x_o,y_o) - \Eb_{P_1}[u_\beta] \big),\\
D_0^{TX} &= (1-r)\big( \Eb_{F_Y}[v_\beta(t_o,x_o,Y')] - \Eb_{P_0}[v_\beta] \big),\\
D_0^{Y} &= (1-r)\big( \Eb_{F_{TX}}[v_\beta(T,X,y_o)] - \Eb_{P_0}[v_\beta] \big),
\end{aligned}
\label{eq:channels}
\end{equation}
and, for any fixed evaluation point $(t,x,y)$,
\begin{equation}
\IF\big(o;q_r(t,x,y),F\big) = \nabla_\beta q_{r,\beta}(t,x,y)^\top\,\IF\big(o;\beta_\lambda,F\big).
\label{eq:if-qr-chain}
\end{equation}
\end{proposition}

\begin{proof}
Perturb $F_\varepsilon=(1-\varepsilon)F+\varepsilon\delta_o$. The real law is
$P_1(F_\varepsilon)=F_\varepsilon$ itself, so for a fixed integrand $u$,
\begin{equation}
\frac{\partial}{\partial\varepsilon}\Eb_{P_1(F_\varepsilon)}[u]\Big|_{0}
= u(t_o,x_o,y_o)-\Eb_{P_1}[u].
\label{eq:proof-channel1}
\end{equation}
The contrastive law factorizes, $P_0(F_\varepsilon)=F_{TX,\varepsilon}\otimes
F_{Y,\varepsilon}$, with marginals
$F_{TX,\varepsilon}=(1-\varepsilon)F_{TX}+\varepsilon\delta_{(t_o,x_o)}$ and
$F_{Y,\varepsilon}=(1-\varepsilon)F_Y+\varepsilon\delta_{y_o}$. Since
\[
\Eb_{P_0(F_\varepsilon)}[v]=\iint v(t,x,y')\,\mathrm dF_{TX,\varepsilon}(t,x)\,\mathrm dF_{Y,\varepsilon}(y')
\]
is a product of two affine functions of $\varepsilon$ paired through the double
integral, the product rule gives one term per factor,
\begin{equation}
\frac{\partial}{\partial\varepsilon}\Eb_{P_0(F_\varepsilon)}[v]\Big|_{0}
= \big(\Eb_{F_Y}[v(t_o,x_o,Y')]-\Eb_{P_0}[v]\big)
+ \big(\Eb_{F_{TX}}[v(T,X,y_o)]-\Eb_{P_0}[v]\big),
\label{eq:proof-channel0}
\end{equation}
the first bracket from differentiating $F_{TX,\varepsilon}$ at fixed $F_Y$, the
second from differentiating $F_{Y,\varepsilon}$ at fixed $F_{TX}$. The mixing weight
$r$ is a design constant and contributes no derivative.

The estimand solves $G_\lambda(\beta_\lambda(F_\varepsilon);F_\varepsilon)=0$ with
$G_\lambda=\nabla_\beta M_\lambda$. Differentiating this identity in $\varepsilon$ at
$0$, and using $\Hb_\lambda=\nabla_\beta G_\lambda$ nonsingular,
\[
\Hb_\lambda\,\frac{\partial}{\partial\varepsilon}\beta_\lambda(F_\varepsilon)\Big|_0
+ \frac{\partial}{\partial\varepsilon}G_\lambda(\beta;F_\varepsilon)\Big|_0 = 0
\quad\Longrightarrow\quad
\IF(o;\beta_\lambda,F) = -\Hb_\lambda^{-1}\,\frac{\partial}{\partial\varepsilon}G_\lambda(\beta;F_\varepsilon)\Big|_0.
\]
Since $G_\lambda(\beta;F_\varepsilon)=r\Eb_{P_1(F_\varepsilon)}[u_\beta]+(1-r)\Eb_{P_0(F_\varepsilon)}[v_\beta]+\lambda\nabla\mathcal{J}(\beta)$
and the penalty does not depend on $F$, applying~\eqref{eq:proof-channel1} to the
$P_1$ term with $u=u_\beta$ and~\eqref{eq:proof-channel0} to the $P_0$ term with
$v=v_\beta$ yields exactly $D_1+D_0^{TX}+D_0^{Y}$, hence~\eqref{eq:master}. Finally
\eqref{eq:if-qr-chain} is the chain rule applied to $\beta\mapsto q_{r,\beta}(t,x,y)$
at fixed $(t,x,y)$.
\end{proof}

Two features of~\eqref{eq:master} matter downstream. The channels $D_0^{TX}$ and
$D_0^{Y}$ are the price of manufacturing the contrastive class from the data: they
are absent from every textbook sandwich formula, which conditions on a fixed training
law, and they vanish only if the contrast distribution is fixed by design, as in
noise-contrastive estimation \citep{gutmann2010nce}, in which case~\eqref{eq:master}
reduces to the textbook formula. Second, no kernel and no bandwidth appear: $q_r$ is
evaluated at a fixed point through the smooth model $q_{r,\beta}$, not through a data
point mass, so the correction term needs no kernel in $T$.

Substituting~\eqref{eq:if-qr-chain} into the correction term
of~\eqref{eq:if-mcd-native} collapses the $X$-expectation into a fixed $k$-vector.

\begin{corollary}[Representer collapse]
\label{cor:collapse}
With
$a_\lambda(t,y):=\Eb_X\!\big[\nabla_\beta q_{r,\beta}(t,X,y)/(1-q_{r,\beta}(t,X,y))^2\big]$,
\begin{equation}
\Eb_X\!\left[\frac{\IF\big(o;q_r(t,X,y),F\big)}{\big(1-q_r(t,X,y)\big)^2}\right]
= a_\lambda(t,y)^\top\,\IF\big(o;\beta_\lambda,F\big),
\label{eq:collapse}
\end{equation}
so the correction term of~\eqref{eq:if-mcd-native} is the scalar
$f_Y(y)\,c\,a_\lambda(t,y)^\top\IF(o;\beta_\lambda,F)$: one linear solve against
$\Hb_\lambda$ and one inner product per evaluation point $(t,y)$.
\end{corollary}

It remains to make $\Hb_\lambda$, the three channels and $a_\lambda$ explicit. We do
this once, for the logistic model, and then show that the MLP and CatBoost reduce to
it.

\subsection{Logistic regression: classical and ridge}
\label{sec:if-qr-logistic}

Let $q_{r,\beta}(t,x,y)=\sigma(\beta^\top\phi(t,x,y))$ with $\sigma(u)=1/(1+e^{-u})$
and a fixed feature map $\phi:\Tset\times\Xset\times\mathbb{R}\to\mathbb{R}^k$, fit by
cross-entropy $\ell(z,q)=-z\log q-(1-z)\log(1-q)$, with penalty
$\mathcal{J}(\beta)=\|\beta\|_2^2$ for ridge and $\lambda=0$ for the classical
unpenalized fit. Abbreviate $q:=q_{r,\beta}$ and $\phi:=\phi(t,x,y)$.

\paragraph{Gradients and Hessian.}
With $s:=\beta^\top\phi$ and $q=\sigma(s)$, using $\sigma'=\sigma(1-\sigma)$,
\[
\nabla_\beta\ell(z,q)
= \frac{\partial\ell}{\partial q}\cdot\frac{\partial q}{\partial s}\cdot\nabla_\beta s
= \Big(-\frac{z}{q}+\frac{1-z}{1-q}\Big)\cdot q(1-q)\cdot\phi
= \big(q-z\big)\,\phi,
\]
the classical cancellation of the sigmoid against the cross-entropy. Hence
\begin{equation}
u_\beta=(q-1)\phi,\qquad v_\beta=q\,\phi,
\qquad
\nabla_\beta u_\beta=\nabla_\beta v_\beta = q(1-q)\,\phi\phi^\top,
\label{eq:logistic-grads}
\end{equation}
the last equality from $\nabla_\beta q=q(1-q)\phi$. Both labels contribute the same
second derivative, so the two expectations in $\Hb_\lambda$ merge into a single one
under the augmented mixture,
\begin{equation}
\Hb_\lambda = \Eb_{P_r}\big[q(1-q)\,\phi\phi^\top\big] + 2\lambda I_k,
\label{eq:logistic-hessian}
\end{equation}
positive definite for $\lambda>0$, and positive definite for $\lambda=0$ as soon as
$\phi$ exhibits no perfect collinearity on the support. In particular $\Hb_\lambda$ is
estimated by one average over the augmented sample $\mathcal{D}_Z$.

\paragraph{Resolving the centering terms.}
The population first-order condition of~\eqref{eq:pop-criterion} reads
\begin{equation}
r\,\Eb_{P_1}[u_\beta]+(1-r)\,\Eb_{P_0}[v_\beta] = -2\lambda\beta_\lambda,
\label{eq:logistic-foc}
\end{equation}
which eliminates the centering terms appearing in~\eqref{eq:channels}. Summing the
three channels,
\[
D_1+D_0^{TX}+D_0^{Y}
= r\,u_\beta(o) + (1-r)\Eb_{F_Y}[v_\beta(t_o,x_o,Y')] + (1-r)\Eb_{F_{TX}}[v_\beta(T,X,y_o)]
- \Big(r\Eb_{P_1}[u_\beta]+2(1-r)\Eb_{P_0}[v_\beta]\Big),
\]
the last bracket collecting one centering term from $D_1$ and two identical ones from
$D_0^{TX}$ and $D_0^{Y}$. Using~\eqref{eq:logistic-foc},
\[
-\Big(r\Eb_{P_1}[u_\beta]+2(1-r)\Eb_{P_0}[v_\beta]\Big)
= 2\lambda\beta_\lambda-(1-r)\Eb_{P_0}[v_\beta],
\]
and absorbing the leftover $-(1-r)\Eb_{P_0}[v_\beta]$ back into either contrastive
channel's centering (they are identical) leaves the compact form below.

\begin{keyresult}
For logistic regression with ridge penalty $\lambda\ge0$, where $\lambda=0$ is the
classical fit,
\begin{equation}
\IF\big(o;\beta_\lambda,F\big)
= -\,\Hb_\lambda^{-1}\Big\{
\underbrace{r\big(q(o)-1\big)\phi(o)}_{\text{real channel}}
+\underbrace{(1-r)\,\Eb_{F_Y}\!\big[q\phi\big](t_o,x_o,\cdot)}_{\text{contrastive }(t,x)\text{-channel}}
+\underbrace{(1-r)\,\Eb_{F_{TX}}\!\big[q\phi\big](\cdot,\cdot,y_o)}_{\text{contrastive }y\text{-channel}}
+2\lambda\beta_\lambda\Big\},
\label{eq:if-beta-logistic}
\end{equation}
with
$\Eb_{F_Y}[q\phi](t_o,x_o,\cdot)=\Eb_{Y'\sim f_Y}\big[q_{r,\beta}(t_o,x_o,Y')\,\phi(t_o,x_o,Y')\big]$
and
$\Eb_{F_{TX}}[q\phi](\cdot,\cdot,y_o)=\Eb_{(T,X)\sim f_{TX}}\big[q_{r,\beta}(T,X,y_o)\,\phi(T,X,y_o)\big]$.
The representer of Corollary~\ref{cor:collapse} collapses to the odds-weighted feature
average
\begin{equation}
a_\lambda(t,y)
= \Eb_X\!\left[\frac{\nabla_\beta q_{r,\beta}(t,X,y)}{(1-q_{r,\beta}(t,X,y))^2}\right]
= \Eb_X\!\left[\frac{q(1-q)}{(1-q)^2}\,\phi\right]
= \Eb_X\big[g(t,X,y)\,\phi(t,X,y)\big],
\label{eq:a-logistic}
\end{equation}
using $q(1-q)/(1-q)^2=q/(1-q)=g$. The correction term of~\eqref{eq:if-mcd-native} is
then $f_Y(y)\,c\,a_\lambda(t,y)^\top\IF(o;\beta_\lambda,F)$.
\end{keyresult}

Setting $\lambda=0$ recovers the classical logistic case: the $2\lambda\beta_\lambda$
term drops and $\Hb_0=\Eb_{P_r}[q(1-q)\phi\phi^\top]$. Dropping in addition the two
contrastive channels, which is what fixing the contrast law by design would do,
leaves $-\Hb_0^{-1}(q(o)-z_o)\phi(o)$, the textbook logistic influence function: a
useful check on~\eqref{eq:if-beta-logistic}.

\paragraph{Estimator.}
Split the augmented sample into an estimation fold and an evaluation fold. On the
estimation fold, fit $\hat\beta_\lambda$ by penalized cross-entropy and form
\[
\hat\Hb_\lambda = \frac1N\sum_{w\in\mathcal{D}_Z} \hat q(w)\big(1-\hat q(w)\big)\phi(w)\phi(w)^\top + 2\lambda I_k,
\qquad \hat q(w)=\sigma\big(\hat\beta_\lambda^\top\phi(w)\big).
\]
Estimate the two contrastive-channel expectations by empirical marginal averages,
using the observed $Y_j$ for the $f_Y$-average and the observed pairs $(T_j,X_j)$ for
the $f_{TX}$-average,
\[
\widehat{\Eb_{F_Y}[q\phi]}(t_o,x_o,\cdot)=\frac1n\sum_j \hat q(t_o,x_o,Y_j)\,\phi(t_o,x_o,Y_j),
\qquad
\widehat{\Eb_{F_{TX}}[q\phi]}(\cdot,\cdot,y_o)=\frac1n\sum_j \hat q(T_j,X_j,y_o)\,\phi(T_j,X_j,y_o),
\]
and set $\hat a_\lambda(t,y)=\frac1n\sum_j \hat g(t,X_j,y)\,\phi(t,X_j,y)$ with
$\hat g=\hat q/(1-\hat q)$. On the evaluation fold, assemble
$\widehat\IF(O_i;\beta_\lambda,\hat F)$ from~\eqref{eq:if-beta-logistic} for each $i$
and plug into the empirical version of~\eqref{eq:if-mcd-native},
\begin{equation}
\hat p_t^{\,\mathrm{MCD}}(y)
= \frac1n\sum_{i} \Big\{
\hat f_Y(y)\,\hat c\Big[\hat g(t,X_i,y) - \tfrac1n\textstyle\sum_j \hat g(t,X_j,y)
+ \hat a_\lambda(t,y)^\top\widehat\IF(O_i;\beta_\lambda,\hat F)\Big]
+ \hat m(t,y)\big(K_{h_Y}(Y_i-y)-\hat f_Y(y)\big)\Big\},
\label{eq:mcd-estimator}
\end{equation}
with $\hat m(t,y)=\hat c\,\frac1n\sum_j\hat g(t,X_j,y)$. The only bandwidth in the
whole estimator is $h_Y$, carried by the marginal correction; there is no kernel in
$T$ and no propensity score anywhere.

\paragraph{The role of $\lambda$.}
At fixed $\lambda>0$ the estimand is the ridge population minimizer $\beta_\lambda(F)$
and~\eqref{eq:if-beta-logistic} is exact for it; the $2\lambda I_k$ term inside
$\Hb_\lambda$ keeps the influence bounded, the local robustness of penalized
M-estimators \citep{avella2017}, and the regularization bias of order $\lambda$ enters
$p_t(y)$ only through the doubly robust remainder, multiplied by the error in the
other nuisance. At $\lambda_n=o(n^{-1/2})$ the penalty is first-order negligible and
\eqref{eq:if-beta-logistic} reduces to the classical influence function. Both regimes
leave the closed form usable as written.

\subsection{Multilayer perceptron and CatBoost through a frozen representation}
\label{sec:if-qr-frozen}

Neither the MLP nor CatBoost satisfies Assumption~\ref{as:smooth}, for the reasons
given in Section~\ref{sec:if-qr-regularity}. Rather than approximate a derivative that
does not exist, we make the estimand smooth by construction, exploiting the fact that
both classifiers already end in a logistic layer.

\paragraph{The two-stage construction.}
Both models factor as a representation followed by a linear-logistic head. For a
network $q=\sigma(l_\theta)$ with $L$ layers, let $\phi(t,x,y)\in\mathbb{R}^k$ denote
the penultimate-layer activation, so that the logit is $l_\theta=\beta^\top\phi$ with
$\beta$ the final-layer weights. For a gradient-boosted ensemble with $M$ trees, let
$\phi(t,x,y)\in\{0,1\}^k$ be the concatenated one-hot leaf-indicator vector,
$k=\sum_{m=1}^M L_m$ with $L_m$ the number of leaves of tree $m$: the entry for leaf
$\ell$ of tree $m$ equals one if $(t,x,y)$ falls into that leaf, zero otherwise. Since
each tree returns the value of the leaf reached and the ensemble sums these values, the
boosted score is exactly linear in $\phi$,
\[
F_M(t,x,y) = \sum_{m=1}^M \nu\,h_m(t,x,y) = \beta^\top\phi(t,x,y),
\]
with $\beta$ collecting the shrunken leaf values. Both classifiers are therefore
already of the form $q_{r,\beta}=\sigma(\beta^\top\phi)$ of
Section~\ref{sec:if-qr-logistic}; what disqualifies them is that $\phi$ is itself
learned from $F$, and learned nonsmoothly.

The construction removes that dependence by cross-fitting. Partition the augmented
sample into folds. On fold $\mathcal{A}$, train the MLP or CatBoost without
restriction and extract the representation $\hat\phi$, that is, the penultimate layer
or the leaf-indicator map. Freeze it: $\hat\phi$ is now a fixed, known function, no
longer an estimated object. On fold $\mathcal{B}$, discard the head learned on
$\mathcal{A}$ and refit a ridge-penalized logistic head
\[
\hat\beta_\lambda = \arg\min_\beta\; \frac{1}{|\mathcal{B}|}\sum_{w\in\mathcal{B}} \ell\big(z_w,\sigma(\beta^\top\hat\phi(w))\big) + \lambda\|\beta\|_2^2 ,
\]
and use $\hat q_r=\sigma(\hat\beta_\lambda^\top\hat\phi)$ as the MCD nuisance. Swap the
folds and average, as usual.

\begin{proposition}[Closed-form influence function under a frozen representation]
\label{prop:frozen}
Conditionally on $\hat\phi$, the head $q_{r,\beta}=\sigma(\beta^\top\hat\phi)$ is a
ridge-penalized logistic model with a fixed feature map, and it satisfies
Assumption~\ref{as:smooth} whenever $\lambda>0$. Consequently
Proposition~\ref{prop:master}, the closed form~\eqref{eq:if-beta-logistic}, the
Hessian~\eqref{eq:logistic-hessian}, the representer~\eqref{eq:a-logistic} and the
estimator~\eqref{eq:mcd-estimator} apply verbatim with
$\phi\rightsquigarrow\hat\phi$, for the MLP and for CatBoost alike.
\end{proposition}

\begin{proof}
Conditionally on $\hat\phi$, the criterion on fold $\mathcal{B}$ is the population
criterion~\eqref{eq:pop-criterion} with loss $\ell$ the cross-entropy, model
$\sigma(\beta^\top\hat\phi)$ and penalty $\lambda\|\beta\|_2^2$. It is strictly convex
in $\beta$ for $\lambda>0$, so its minimizer exists, is unique and interior; the loss
and penalty are twice differentiable; and by~\eqref{eq:logistic-hessian} the Hessian
$\Eb_{P_r}[q(1-q)\hat\phi\hat\phi^\top]+2\lambda I_k$ is positive definite, hence
nonsingular. The remaining condition, $q_{r,\beta}$ bounded away from $0$ and $1$,
holds on any bounded parameter region with $\hat\phi$ bounded, which is automatic for
leaf indicators and holds for bounded activations. Assumption~\ref{as:smooth} is
therefore satisfied, and Proposition~\ref{prop:master} applies with feature map
$\hat\phi$.
\end{proof}

Cross-fitting is what licenses treating $\hat\phi$ as fixed while differentiating on
$\mathcal{B}$: the representation is a function of $\mathcal{A}$ only, independent of
$\mathcal{B}$, so the perturbation of $F$ along $\mathcal{B}$ does not move it. This
is the same device that removes empirical-process conditions from every doubly robust
one-step estimator \citep{chernozhukov2018,kennedy2022review}, used here for an
additional purpose: to restore differentiability of the nuisance itself.

\paragraph{What is gained and what is paid.}
What is gained is that all three classifiers share one closed form. The flexibility of
the network or of the boosted ensemble is preserved, in $\hat\phi$: the representation
still learns the interactions, nonlinearities and categorical structure that made those
models attractive in the first place. What is paid is a projection. Differentiating
the head gives the influence function of the model $\sigma(\beta^\top\hat\phi)$, not of
the true $q_r$; equivalently, the correction acts along the tangent space spanned by
$\hat\phi$, the set of directions in which the frozen model can move, and any component
of the true perturbation orthogonal to that span is not corrected. The resulting bias
is second order, a product of the head's misspecification error with the error in the
other nuisance, and is absorbed by the doubly robust remainder of the estimator: the
same regime in which the whole MCD construction already operates. The practical
implication is that $\hat\phi$ should be rich, since the width of the penultimate layer
or the total number of leaves controls the span, and that $k$ is the dimension of the
one linear solve in~\eqref{eq:if-beta-logistic}, which bounds how rich it can be in
practice.

Two implementation points are worth recording. For CatBoost the leaf-indicator map is
sparse, exactly $M$ nonzero entries out of $k$, so $\hat\Hb_\lambda$ is cheap to
accumulate and $\hat a_\lambda$ cheap to evaluate; oblivious trees make the leaf index
a fixed-depth binary code, so $\hat\phi$ is read off directly. For the MLP the ridge
penalty on the head is essential rather than cosmetic, since it is what guarantees the
nonsingularity used in Proposition~\ref{prop:frozen}, and it plays the role that
damping plays in \citet{kohliang2017}, but here as part of the estimand rather than as
a numerical fix. In both cases the probability output should be a consistent estimator
of $q_r$, which cross-entropy training delivers, being a proper composite loss; boosted
probabilities are known to be distorted toward a sigmoidal shape in finite samples
\citep{niculescu2005predicting}, and the refitted logistic head performs exactly the
recalibration that would otherwise be added post hoc.

\begin{remark}[Fallback: the representer route]
If one refuses to freeze the representation and insists on debiasing the full,
end-to-end trained classifier, the nuisance-level influence function must be abandoned
and the correction reorganized. The correction term of~\eqref{eq:if-mcd-native} is the
continuous linear functional
$\dot q\mapsto f_Y(y)\,c\,\Eb_X[\dot q(t,X,y)/(1-q_r(t,X,y))^2]$ evaluated at
$\dot q=\IF(o;q_r,F)$, and by Riesz representation this evaluation equals
$\Eb[\alpha(O)\,s(O;o)]$ for a representer $\alpha$, which the equivalence between the
MCD-native and classical corrections identifies as the localized inverse generalized
propensity score $\alpha(t;O)=K_{h_T}(T-t)/\pi(t\mid X)$. Automatic debiased machine
learning estimates $\alpha$ by minimizing a Riesz loss over a flexible class
\citep{chernozhukov2022autodml,chernozhukov2022riesznet,kato2026riesz}, with the
kernel-localized continuous-treatment version of \citet{colangelo2020}. We do not take
this route. Although it estimates $\alpha$ by regression rather than by conditional
density estimation, it reintroduces the propensity score as the object being learned
and, since $\alpha$ is of order $h_T^{-1}$, it reintroduces the treatment dimension
through the variance of a localized regression, which is precisely what MCD was
designed to avoid. The frozen-representation route above keeps the correction inside
the classifier and pays instead a second-order projection bias.
\end{remark}

\paragraph{Summary.}
For all three classifiers the influence function of the nuisance is the closed
form~\eqref{eq:if-beta-logistic}, with representer~\eqref{eq:a-logistic}, plugged into
the estimator~\eqref{eq:mcd-estimator}. Classical and ridge logistic regression satisfy
the required regularity directly. The MLP and CatBoost satisfy it after the
representation is frozen on an independent fold and the logistic head refitted, which
costs a second-order projection bias absorbed by the doubly robust remainder and buys a
correction expressed entirely through the classifier's own parameters, with no
generalized propensity score and no kernel in the treatment coordinate.

\bibliographystyle{plainnat}
\bibliography{references}

% New BibTeX keys to add to references.bib (the rest already exist in the paper):
% prokhorenkova2018catboost, friedman2001greedy, niculescu2005predicting.

\end{document}